\documentclass[10pt]{article}

\usepackage[utf8]{inputenc}
\usepackage[T1]{fontenc}
\usepackage[spanish]{babel}
\usepackage[a4paper]{geometry}
\usepackage{amsmath}
\usepackage{multicol}
\usepackage{enumitem}
\usepackage{fancyhdr}
\usepackage{titlesec}

% Configuración de página
\pagestyle{fancy}
\fancyhf{}
\fancyhead[L]{Academia Preuniversitaria - CTA}
\fancyhead[R]{\today}
\renewcommand{\headrulewidth}{0.4pt}

% Configuración de títulos
\titleformat{\section}[block]{\large\bfseries\centering}{\thesection}{1em}{}

\setlength{\columnsep}{1cm}
\setlength{\parskip}{0.5em}
\setlength{\parindent}{0em}

% Alternativas
\newlist{choices}{enumerate}{1}
\setlist[choices]{label=\alph*), itemsep=0.15em}

\begin{document}
	
	\begin{center}
		\Large\bfseries Hoja de Ejercicios de CTA
		
		\normalsize Materia, Energía y el Universo
	\end{center}
	
	\begin{multicols}{2}
		
		\section*{Ejercicios}
		
		\begin{enumerate}[label=\textbf{\arabic*.}]
			
			\item La materia se caracteriza porque:
			\begin{choices}
				\item Solo posee energía
				\item Tiene masa y ocupa espacio
				\item Carece de volumen
				\item Es invisible
				\item Solo existe en estado sólido
			\end{choices}
			
			\vspace{0.3cm}
			
			\item ¿Cuál es una propiedad general de la materia?
			\begin{choices}
				\item Elasticidad
				\item Ductilidad
				\item Solubilidad
				\item Masa
				\item Magnetismo
			\end{choices}
			
			\vspace{0.3cm}
			
			\item La densidad se calcula mediante:
			\begin{choices}
				\item $\rho = m \cdot V$
				\item $\rho = \dfrac{V}{m}$
				\item $\rho = \dfrac{m}{V}$
				\item $\rho = mg$
				\item $\rho = \dfrac{F}{A}$
			\end{choices}
			
			\vspace{0.3cm}
			
			\item El estado de la materia con volumen fijo y forma variable es:
			\begin{choices}
				\item Plasma
				\item Sólido
				\item Líquido
				\item Gaseoso
				\item Vapor
			\end{choices}
			
			\vspace{0.3cm}
			
			\item El cambio de estado de sólido a líquido se denomina:
			\begin{choices}
				\item Condensación
				\item Solidificación
				\item Vaporización
				\item Fusión
				\item Sublimación
			\end{choices}
			
			\vspace{0.3cm}
			
			\item La energía es:
			\begin{choices}
				\item La capacidad de ocupar espacio
				\item La resistencia al movimiento
				\item La capacidad de producir cambios
				\item Una propiedad exclusiva de los sólidos
				\item La fuerza gravitacional
			\end{choices}
			
			\vspace{0.3cm}
			
			\item Según la Ley de Conservación de la Energía:
			\begin{choices}
				\item La energía desaparece
				\item La energía solo se crea
				\item La energía se transforma
				\item La energía aumenta indefinidamente
				\item La energía no produce cambios
			\end{choices}
			
			\vspace{0.3cm}
			
			\item La energía asociada al movimiento es:
			\begin{choices}
				\item Nuclear
				\item Potencial
				\item Térmica
				\item Cinética
				\item Química
			\end{choices}
			
			\vspace{0.3cm}
			
			\item ¿Cuál es una fuente de energía renovable?
			\begin{choices}
				\item Petróleo
				\item Carbón
				\item Gas natural
				\item Energía solar
				\item Uranio
			\end{choices}
			
			\vspace{0.3cm}
			
			\item El centro del Sistema Solar es:
			\begin{choices}
				\item La Tierra
				\item Marte
				\item Júpiter
				\item La Luna
				\item El Sol
			\end{choices}
			
			\vspace{0.3cm}
			
			\item El planeta más grande del Sistema Solar es:
			\begin{choices}
				\item Saturno
				\item Marte
				\item Júpiter
				\item Venus
				\item Neptuno
			\end{choices}
			
			\vspace{0.3cm}
			
			\item El movimiento de rotación de la Tierra produce:
			\begin{choices}
				\item Las estaciones
				\item Los eclipses
				\item El día y la noche
				\item Los terremotos
				\item La lluvia
			\end{choices}
			
			\vspace{0.3cm}
			
			\item La capa más externa de la Tierra es:
			\begin{choices}
				\item Núcleo interno
				\item Núcleo externo
				\item Manto
				\item Corteza
				\item Astenósfera
			\end{choices}
			
			\vspace{0.3cm}
			
			\item Un mineral es:
			\begin{choices}
				\item Una mezcla artificial
				\item Una sustancia orgánica
				\item Una sustancia natural e inorgánica
				\item Un líquido viscoso
				\item Un tipo de energía
			\end{choices}
			
			\vspace{0.3cm}
			
			\item El granito es una roca:
			\begin{choices}
				\item Sedimentaria
				\item Metamórfica
				\item Ígnea
				\item Artificial
				\item Orgánica
			\end{choices}
			
			\vspace{0.3cm}
			
			\item La teoría más aceptada sobre el origen del Universo es:
			\begin{choices}
				\item Generación espontánea
				\item Universo estacionario
				\item Big Bang
				\item Geocentrismo
				\item Creacionismo científico
			\end{choices}
			
			\vspace{0.3cm}
			
			\item La expansión del Universo fue demostrada por:
			\begin{choices}
				\item Newton
				\item Einstein
				\item Galileo
				\item Hubble
				\item Copérnico
			\end{choices}
			
			\vspace{0.3cm}
			
			\item Aproximadamente, ¿qué porcentaje del Universo corresponde a materia ordinaria?
			\begin{choices}
				\item 68\%
				\item 50\%
				\item 27\%
				\item 15\%
				\item 5\%
			\end{choices}
			
		\end{enumerate}
		
	\end{multicols}
	
\end{document}